\subsection{PREORDER RELATIONS}

Let \(\mathbb{R}\{x, y\}\) be a relation, \(x\) and \(y\) being distinct letters. If \(\mathbb{R}\) is transitive and if we have \(\mathbb{R}\{x, y\} \Longrightarrow (\mathbb{R}\{x, x\} \text{ and } \mathbb{R}\{y, y\})\) , it does not necessarily follow that \(\mathbb{R}\) is an order relation because the relation

\[
(\mathbf {R} \{x, y \} \text { and } \mathbf {R} \{y, x \})
\]

does not necessarily imply \(x = y\) . \(\mathbb{R}\{x, y\}\) is said to be a preorder relation between \(x\) and \(y\) ; \(\mathbb{R}\{y, x\}\) is then also a preorder relation between \(x\) and \(y\) , called the opposite of the relation \(\mathbb{R}\{x, y\}\) .

For example, let R be the set of subsets of P(E) which are coverings of E (Chapter II, §4, no. 6). The relation ``R is coarser than R'' between elements R, R' of R (Chapter II, §4, no. 6) is transitive and reflexive. But two distinct coverings can be such that each is coarser than the other; for example, this is the case when R' is the union (in P(E)) of R and a subset of E contained in a set of R but not belonging to R.

But in any case the relation \((\mathbb{R}\{x, y\}\) and \(\mathbb{R}\{y, x\})\) is an equivalence relation \(\mathbb{S}\{x, y\}\) with respect to \(x\) and \(y\) . Let \(x', y'\) be letters distinct from \(x, y\) which do not appear in \(\mathbb{R}\) . Then \(\mathbb{R}\{x, y\}\) is compatible (with respect to \(x\) and \(y\) ) with the equivalence relations \(\mathbb{S}\{x, x'\}\) and \(\mathbb{S}\{y, y'\}\) ; in other words (Chapter II, §6, no. 8), the relation

\[
(\mathrm{R} \{x, y \} \text {   and   } \mathrm{S} \{x, x ^ {\prime} \} \text {   and   } \mathrm{S} \{y, y ^ {\prime} \})
\]

implies R \(\{x', y'\}\) .

A preorder relation on a set E is a preorder relation R\{x, y\} such that the relation R\{x, x\} is equivalent to \(x \in E\) . The relation R\{x, y\} then implies `` \(x \in E\) and \(y \in E\) ''.

If \(\mathbf{R}\{x, y\}\) is a preorder relation on a set \(\mathbf{E}\) , then the relation \(\mathbf{S}\{x, y\}\) defined above is an equivalence relation on \(\mathbf{E}\) . Let \(\mathbf{R}'\{\mathbf{X}, \mathbf{Y}\}\) denote the relation

\[
\mathbf {X} \in \mathbf {E} / \mathrm{S} \text { and } \mathbf {Y} \in \mathbf {E} / \mathrm{S} \text { and } (\exists x) (\exists y) (x \in \mathbf {X} \text { and } y \in \mathbf {Y} \text { and } \mathbf {R} \{x, y \}),
\]

that is to say, the relation induced by R on passing to the quotient (with respect to x and y); we saw in Chapter II, §6, no. 3, that it is equivalent to the relation

\[
\mathrm{X} \in \mathrm{E} / \mathrm{S} \text {   and   } \mathrm{Y} \in \mathrm{E} / \mathrm{S} \text {   and   } (\forall x) (\forall y) ((x \in \mathrm{X} \text {   and   } y \in \mathrm{Y}) \implies \mathrm{R} \{x, y \}).
\]

Let us show that \(R'\{X, Y\}\) is an order relation between elements of E/S. The relation \((R'\{X, Y\}\) and \(R'\{Y, Z\})\) is equivalent to

\(\mathbf{X} \in \mathbf{E} / \mathbf{S}\) and \(\mathbf{Y} \in \mathbf{E} / \mathbf{S}\) and \(\mathbf{Z} \in \mathbf{E} / \mathbf{S}\) and \((\forall x)(\forall y)(\forall z)((x \in \mathbf{X}\) and \(y \in \mathbf{Y}\) and \(z \in \mathbf{Z}) \Longrightarrow (\mathbf{R}\{x, y\}\) and \(\mathbf{R}\{y, z\}))\)

(Chapter I, §4, criteria C40, C41). Since R\{x, y\} is transitive and Y ∈ E/S ⟶ Y ≠ φ (Chapter II, §6, no. 2), it follows immediately that R' \{X, Y\} is transitive. Next, (R' \{X, Y\} and R' \{Y, X\}) is equivalent to

\[
\begin{array}{l} \mathrm{X} \in \mathrm{E} / \mathrm{S} \text { and } \mathrm{Y} \in \mathrm{E} / \mathrm{S} \text { and } \\ (\forall x) (\forall y) ((x \in \mathrm{X} \text { and } y \in \mathrm{Y}) \Longrightarrow (\mathrm{R} \{x, y \} \text { and } \mathrm{R} \{y, x \})), \end{array}
\]

and hence equivalent to

\(X \in E/S\) and \(Y \in E/S\) and \((\forall x)(\forall y)((x \in X \text{ and } y \in Y) \Longrightarrow S\{x, y\})\) , and therefore implies

\[
\mathrm{X} \in \mathrm{E} / \mathrm{S} \quad \text { and } \quad \mathrm{Y} \in \mathrm{E} / \mathrm{S} \quad \text { and } \quad \mathrm{X} = \mathrm{Y}.
\]

Moreover, R \(\{x, y\}\) implies R \(\{x, x\}\) and R \(\{y, y\}\) , and hence R' \(\{X, Y\}\) implies each of the relations

\[
\begin{array}{l} \mathrm{X} \in \mathrm{E} / \mathrm{S} \text { and } (\forall x) ((x \in \mathrm{X}) \Longrightarrow \mathrm{R} \{x, x \}), \\ \mathrm{Y} \in \mathrm{E} / \mathrm{S} \text { and } (\forall y) ((y \in \mathrm{Y}) \Longrightarrow \mathrm{R} \{y, y \}), \end{array}
\]

whence \(R'\{X,Y\}\) implies \((R'\{X,X\}\) and \(R'\{Y,Y\})\) . Finally, since \(x \in E\) implies \(R\{x,x\}\) , \(X \in E/S\) implies \(R'\{X,X\}\) , and the proof of our assertion is complete. \(R'\{X,Y\}\) is said to be the order relation associated with \(R\{x,y\}\) .

A preordering on a set E is a correspondence \(\Gamma = (G, E, E)\) with E as source and as target, and such that \((x, y) \in G\) is a preorder relation on E. By abuse of language we refer sometimes to the graph \(G\) of \(\Gamma\) as a preordering on E. For this to be so it is necessary and sufficient that \(\Delta \subset G\) and \(G \circ G \subset G\) (which implies \(G \circ G = G\) ). The equivalence relation S

corresponding to the preorder relation \((x, y) \in G\) then has \(G \cap \overline{G}\) as its graph; the order relation associated with \((x, y) \in G\) has as graph the subset \(G'\) of \((E/S) \times (E/S)\) which corresponds (Chapter II, §6, no. 8) to the image of \(G\) under the canonical mapping of \(E \times E\) onto

\[
(\mathbf {E} \times \mathbf {E}) / (\mathbf {S} \times \mathbf {S}).
\]

Example. * Let A be a ring with an identity element. The relation \((\exists z)(z \in A\) and \(y = zx)\) between two elements \(x\) , \(y\) of A is a preorder relation on A; it is read ``\(x\) is a right divisor of \(y\)'' or ``\(y\) is a left multiple of \(x\)''. *

